% Polska wersja robocza artykułu windowed_qubo — do szybkiego czytania.
% Wiążąca jest wersja angielska (main.tex). Kompilacja: xelatex wersja_pl.tex
\documentclass[11pt,a4paper]{article}
\usepackage{fontspec}
\usepackage{polyglossia}
\setdefaultlanguage{polish}
\usepackage{amsmath,amssymb}
\usepackage{booktabs}
\usepackage{multirow}
\usepackage[margin=2.4cm]{geometry}
\usepackage{parskip}
\usepackage{hyperref}
\hypersetup{colorlinks=true, linkcolor=blue!55!black, urlcolor=blue!55!black}

\title{Skalowalne kwantowe przeszukiwanie sąsiedztwa dla permutacyjnego
  problemu przepływowego przez okienkową dekompozycję QUBO\\
  \large Wersja polska do szybkiego czytania (wiążąca: main.tex, EN)}
\author{Remigiusz Wojewódzki \and Wojciech Bożejko\\
  \small Katedra Sterowania i Obliczeń Kwantowych, WIT, Politechnika Wrocławska}
\date{Stan tłumaczenia: 17.07.2026 --- cel: Computers \& Operations Research}

\begin{document}
\maketitle

\section*{Abstrakt}
Rozszerzamy framework kwantowych sąsiedztw dla permutacyjnego problemu
przepływowego (PFSP) o trzecią klasę --- \emph{quantum QUBO enhanced} ---
działającą na znacznie większych instancjach, niż pozwalały pierwotne
formulacje na QPU D-Wave. Dotychczasowy framework obejmował cztery
sąsiedztwa klasyczne (Adjacent, Fibonacci, Dynasearch, Motzkin) i cztery
kwantowe QUBO ograniczone do $n=20$. Klasa enhanced usuwa dwa ograniczenia:
zbędny filtr delty w strukturach liniowych (Adjacent/Fibonacci), co
rozszerza stosowalność QPU do $n\le200$ i $n\le500$, oraz wprowadza
okienkową dekompozycję QUBO dla struktur kwadratowych (Dynasearch/Motzkin),
przesuwając wykonalność do $n\le50$. Oba frameworki metaheurystyczne (ILS
z dywersyfikacją tabu przez elitarną pulę BackTrackJumpList oraz SA
aplikujące akceptację Metropolisa do najlepszego sąsiada z pełnego skanu,
z adaptacyjnym podgrzewaniem) używają akceleratora własności blokowych,
tnącego liczbę zmiennych QUBO o 30--60\%. Dwanaście sąsiedztw oceniamy na
benchmarkach Taillarda przy sześciu budżetach (100--10\,000 ms). Pomiary
QPU na Advantage\_system4.1 potwierdzają skalowanie: średni RPD
fibonacciego enhanced spada monotonicznie z 31{,}7\% ($n=20$) do 15{,}0\%
($n=200$), a luka do klasycznego fibonacciego zwęża się z 20{,}4~pp
($n=20$) do 4{,}6~pp ($n=500$), gdzie wariant kwantowy dorównuje
klasycznemu adjacent. Okienkowanie jest niemal darmowe tam, gdzie
interakcje swapów są lokalne (dynasearch), a kosztuje ok.~4~pp ---
rosnąco z budżetem --- tam, gdzie są globalne (motzkin).

\paragraph{Highlights.} Klasa enhanced łamie barierę $n=20$ na QPU D-Wave;
usunięcie filtra delty rozszerza Adjacent/Fibonacci do $n\le200$/$n\le500$;
okienkowa dekompozycja umożliwia Dynasearch i Motzkin do $n=50$;
akcelerator Smutnickiego tnie zmienne QUBO o 30--60\%; luka
kwantowo-klasyczna zwęża się z 20{,}4~pp do 4{,}6~pp; okienkowanie prawie
darmowe dla struktur lokalnych, strata narasta dla globalnych.

\section{Wstęp}
PFSP szuka permutacji $n$ zadań na $m$ maszynach szeregowo minimalizującej
makespan $C_{\max}$. Problem jest NP-trudny dla $m\ge3$
(Garey--Johnson--Sethi 1976), a dla dużych instancji najlepsze rozwiązania
dają metaheurystyki oparte na sąsiedztwach (monografia Bożejko 2010).
Wybór struktury sąsiedztwa --- zbioru rozwiązań osiągalnych jednym
dopuszczalnym ruchem (van Laarhoven i in.\ 1992) --- jest głównym motorem
jakości rozwiązań.

\paragraph{Klasyczna linia flow shopu.} Johnson (1954) rozwiązuje
przypadek dwóch maszyn wielomianowo; heurystyka NEH (1983) pozostaje
standardowym punktem startu. Linia metaheurystyczna przeszła przez SA
(Osman--Potts 1989) i algorytmy genetyczne (Reeves 1995), ale jej
kręgosłupem jest tabu search: implementacja Taillarda (1990), blokowy
fast tabu Nowickiego--Smutnickiego (1996a) i przyspieszony następca
Grabowskiego--Wodeckiego (2004) kolejno cięły czas i dystans do optimum
na benchmarku Taillarda (1993). Pole skonsolidowały przegląd
Ruiza--Maroto (2005) i iterated greedy Ruiza--Stützlego (2007), który
pokazał, że prosta pętla destrukcja--rekonstrukcja dorównuje bardziej
wyszukanym maszynom.

\paragraph{Linia otoczeń.} W każdym mocnym algorytmie PFSP powtarzają się
dwa składniki: struktura ruchu i własności ścieżki krytycznej, które ją
przycinają. Dynasearch (Congram--Potts--van de Velde 2002) pokazał, że
wykładniczo duży zbiór parami niezależnych ruchów da się przeszukać
dokładnie programowaniem dynamicznym w czasie wielomianowym ---
jednostką przeszukiwania staje się ruch złożony. Teorię multimoves
rozwinęli Bożejko--Wodecki (2007), a paradygmat rozrósł się w very
large-scale neighborhood search (Ahuja i in.\ 2002). Własności blokowe
(Nowicki--Smutnicki, Smutnicki 1998, Grabowski--Wodecki) atakują to samo
wąskie gardło z drugiej strony: mówią, których kandydatów można pominąć
hurtem. Nasze sąsiedztwa Fibonacci i Motzkin (ICAISC 2026; EXIT 2025)
kontynuują linię ruchów złożonych: zbiory ruchów zliczane klasycznymi
ciągami, przeszukiwane przez DP i --- co ten artykuł wykorzystuje ---
wyrażalne jako QUBO.

\paragraph{Kwanty w harmonogramowaniu.} Weszły od strony job shopu:
Venturelli i in.\ (całe instancje jako QUBO na annealerze), Kurowski
i in.\ (hybrydy annealera z heurystykami), Amaro i in.\ (obwody
wariacyjne). Wszystkie kodują cały harmonogram naraz, co ogranicza je do
rozmiarów zabawkowych. Nasza alternatywa: QPU \emph{wewnątrz}
metaheurystyki zamiast w jej miejsce --- QUBO staje się sąsiedztwo, nie
harmonogram.

W pracy BPASTS sformułowaliśmy cztery klasyczne sąsiedztwa jako QUBO
i wykonaliśmy je na QPU D-Wave; zakres ograniczał się do $n=20$
(pojemność embeddingu ok.\ 170--180 zmiennych logicznych dla gęstych
grafów na topologii Pegasus P16 z 5627 aktywnymi kubitami). Ten artykuł
usuwa oba blokery. Pierwszym jest zbędny filtr delty: obejście pod
klasyczny symulator, szkodliwe na QPU, bo odrzuca informację o strukturze
kar, podczas gdy QPU sam przypisuje $x_k=0$ ruchom niepoprawiającym;
dla Adjacent i Fibonacciego ($K=n-1$ niezależnie od filtra) usunięcie nic
nie kosztuje. Drugim jest sufit zmiennych $O(n^2)$ dla Dynasearcha
i Motzkina: wprowadzamy dekompozycję okienkową --- nakładające się okna
rozmiaru $w\le19$, QUBO na okno, scalanie zachłanne. Razem daje to
dwanaście sąsiedztw w trzech klasach, po raz pierwszy pokrywających
pełen zakres benchmarku Taillarda.

\paragraph{Kontrybucje.} Praca wnosi okienkową dekompozycję QUBO
(QPU do $n=50$, automatyczny dobór okna z pojemności sprzętu), formulacje
bez filtra delty (Adjacent $n\le200$, Fibonacci $n\le500$), akcelerator
blokowy Smutnickiego we wszystkich dwunastu sąsiedztwach (30--60\% mniej
zmiennych), BackTrackJumpList zamiast restartów losowych w ILS oraz
pierwsze porównanie trzech klas na $n\in\{20,50,100,200,500\}$ przy
identycznych frameworkach i budżetach.

\section{Definicja problemu}
Materiał standardowy, przytaczany dla samodzielności tekstu, w notacji
klasycznej pracy Nowickiego--Smutnickiego (1996a).

\paragraph{PFSP.} Zbiór zadań $J=\{1..n\}$ i maszyn $M=\{1..m\}$;
zadanie $j$ to ciąg operacji $O_{1j}..O_{mj}$, operacja $O_{ij}$ na
maszynie $i$ trwa $p_{ij}>0$ bez przerwań; każda maszyna przetwarza
w tej samej kolejności, więc porządek to permutacja $\pi\in\Pi$.
Czasy zakończenia spełniają
$C_{i,\pi(j)} = \max\{C_{i-1,\pi(j)},\,C_{i,\pi(j-1)}\} + p_{i,\pi(j)}$
z zerowymi warunkami brzegowymi; $C_{\max}(\pi)=C_{m,\pi(n)}$, szukamy
$\pi^*$ minimalizującej $C_{\max}$. Równoważnie makespan to długość
najdłuższej (krytycznej) ścieżki w grafie siatkowym:
\[ C_{\max}(\pi) = \max_{1=j_0\le j_1\le\cdots\le j_m=n}\;
   \sum_{i=1}^{m}\sum_{j=j_{i-1}}^{j_i} p_{i,\pi(j)}, \]
a maksymalne podciągi ścieżki krytycznej na jednej maszynie tworzą
\emph{bloki} --- strukturę, którą akcelerator zamienia w filtr kandydatów.
Pełny opis: Nowicki--Smutnicki 1996a i monografia Smutnickiego 2012;
relacje między permutacjami mają też ujęcie algebraiczne przez funkcje
dodatnio określone na grupach permutacji (Bożejko--Bożejko 2015).

\paragraph{Dekompozycja Head--Tail} (Wodecki--Bożejko 2004). Delta ruchu
w $O(1)$ z macierzy $H[i][j]$ (najwcześniejsze zakończenia od lewej)
i $T[i][j]$ (wkłady od prawej), budowanych w $O(mn)$ raz na iterację.

\paragraph{Metaheurystyki.} Samo sąsiedztwo to jeszcze nie algorytm ---
potrzebna jest warstwa sterująca akceptująca ruchy i wyciągająca
trajektorię z optimów lokalnych. Żeby nie zaburzała porównania, każde
sąsiedztwo działa w tych samych dwóch frameworkach z identycznymi
parametrami --- jedyną zmienną jest sąsiedztwo.

\emph{ILS} (Lourenço i in.\ 2003) działa z listą tabu (Glover 1989)
o tenurze $\tau=10$ i kryterium aspiracji przyjmującym ruch tabu, jeśli
ściśle poprawia globalne minimum --- w duchu fast taboo search
Nowickiego--Smutnickiego (1996b).

\emph{SA} (Kirkpatrick 1983) działa z $T_0=1000$, $\alpha=0.99$, podłogą
$T_{\min}=100$ i dywersyfikacją reheat-kick. \textbf{Nasz wariant różni
się od kanonicznego SA kryterium eksploracji} (komponent frameworku
Franzina--Stützle 2019): kanoniczne SA losuje sąsiada; Ishibuchi i in.\
(1995) --- właśnie dla flow shopu --- poddawali Metropolisowi najlepszego
z losowej próbki sąsiadów, co uodparnia SA na wybór harmonogramu
chłodzenia. My doprowadzamy ten mechanizm do granicy wyczerpującej:
test dostaje najlepszy ruch \emph{całego} otoczenia ze wspólnego
oraculum. Granica nie jest niewinna: po odrzuceniu losowej próbki losuje
się świeżą, a deterministyczny pełny skan re-proponuje ten sam ruch przy
coraz niższym $T$ --- to patologia zamarzania, którą diagnozuje
sekcja~2.2. Konstrukcję wymusza porównanie: wywołanie QPU zwraca
najlepszy ruch złożony, nie losowy. Ewaluacja pełnego otoczenia pod
akceptacją Metropolisa to ugruntowany reżim w sprzęcie QUBO
(parallel-trial Digital Annealera --- Aramon i in.\ 2019) i w rodzinie
rejection-free (Rosenthal i in.\ 2021; następca losowany z wag zamiast
argmax). Według naszej wiedzy łańcuch ,,wyczerpujący argmax +
Metropolis'' nie został w literaturze nazwany ani przeanalizowany ---
dokumentujemy go empirycznie.

\subsection{BackTrackJumpList}
Restart losowy wyrzuca wszystko, czego szukanie się nauczyło; BTJL to
zachowuje. Pula $k=10$ najlepszych różnych permutacji; przy konflikcie
tabu bez aspiracji algorytm aplikuje perturbację double-bridge (4-opt)
do kolejnej elity (round-robin). Double-bridge tnie $\pi$ w trzech
losowych punktach i skleja $A|C|B|D$ --- ruch spoza otoczenia 2-opt,
więc szukanie naprawdę opuszcza basen; round-robin nie tłucze wciąż
globalnego najlepszego. Pula aktualizowana przy każdej ścisłej poprawie;
przed pierwszą elitą fallbackiem jest restart losowy.

\subsection{SA z reheat-kick}
Wspólne oracula zmuszają krok SA do konsumpcji \emph{najlepszego} ruchu.
Z deterministycznym oraculum szukanie zapada się w minimum lokalnym $X$:
oraculum daje najlepszego sąsiada $X'$ ($\Delta>0$), Metropolis czasem
przyjmie, z $X'$ oraculum wraca do $X$ ($\Delta<0$, zawsze). Trajektoria
wpada w dwucykl $X\leftrightarrow X'$, a podgrzewanie samo go nie łamie.
Rozwiązanie sprzęga reheat z perturbacją w przestrzeni rozwiązań: przy
stagnacji $s=100$ ms bez poprawy (1) $T\leftarrow T_0$, (2) bieżące
rozwiązanie zastępuje double-bridge kolejnej elity, (3) reset licznika.
Podłoga $T_{\min}=100$ utrzymuje ruchliwość pod górę. Dwucykl znika,
a SA zbiega z trajektoriami zależnymi od seeda.

\section{Taksonomia sąsiedztw}
Definicje klasyczne są znane; przypominamy je, bo konstrukcje QUBO
budują wprost na ich zbiorach ruchów.

\subsection{Klasyczne}
\textbf{Adjacent} ($N_{\mathrm{adj}}$): ruch zamienia zadania na
pozycjach $i$ oraz $i{+}1$, więc $|N_{\mathrm{adj}}|=n-1$; w literaturze
API --- adjacent pairwise interchange (hierarchia uogólnionych wymian:
Della Croce 1995).
\textbf{Fibonacci} ($N_{\mathrm{fib}}$): ruch aplikuje jednocześnie
zbiór niekolidujących swapów sąsiednich ($S$ poprawny, gdy $i\in S
\Rightarrow i{+}1\notin S$); liczba podzbiorów $=F_{n+1}$, stąd nazwa
(wprowadzone w ICAISC 2026); optymalny podzbiór wybiera DP w $O(n)$.
\textbf{Dynasearch} ($N_{\mathrm{dyn}}$): swapy końców odcinków dowolnej
długości; DP (Congram 2002) wybiera podzbiór par przedziałowo
nienakładających się, łącznie minimalizując $C_{\max}$; $K=O(n^2)$,
DP w $O(n^3)$, akceleracja filtrem NPI.
\textbf{Motzkin} ($N_{\mathrm{motz}}$): te same swapy końców, ale
dopuszczalne podzbiory wg kombinatoryki ścieżek Motzkina (Motzkin 1948;
SOCO 2026; EXIT 2025): bez krzyżowań i wspólnych końców, ścisłe
zagnieżdżenie dozwolone; liczba podzbiorów $=M_n$ ($M_5=21$,
$M_{10}=4862$, $M_{20}\approx1{,}7\cdot10^{10}$); DP w $O(n^3)$.

\subsection{Akceleratory blokowe}
Pozycja $u$ jest granicą bloku, gdy $H[m][u]=H[m-1][u]+p_{m,\pi(u+1)}$;
zbiór granic $\mathcal U(\pi)$ liczony w $O(mn)$. Twierdzenie
(Smutnicki 1998): swap $(i,j)$ może ściśle poprawić $C_{\max}$ tylko gdy
$\mathcal U(\pi)\cap[i,j]\neq\emptyset$. Adjacent/Fibonacci: z $n-1$
kandydatów zostaje $|\mathcal U(\pi)|$ (na tai $m=5$, $n=100$ średnio
4--8 z 99). Dynasearch/Motzkin (NPI): $O(n^2)\to O(n\cdot|\mathcal U|)$;
w QUBO tylko przefiltrowane pary stają się zmiennymi. Na Taillardzie
filtr usuwa 35--65\% par.

\subsection{Kwantowe QUBO --- oryginalne i enhanced}
Oryginalne (za BPASTS): $x_k=1 \Leftrightarrow$ swap $k$ wybrany; QUBO
koduje cel i wykonalność; $n\le20$.
Enhanced: (a) \emph{pełny QUBO bez filtra delty} --- adjacent (gęsty $Q$)
do $n\le200$, fibonacci (tridiagonalny $Q$) do $n\le500$;
(b) \emph{dekompozycja okienkowa} dla $K=O(n^2)$: okno $\ell$ pokrywa
$[s_\ell, s_\ell{+}w)$, $s_\ell=\ell\lfloor w(1{-}\rho)\rfloor$,
$\rho=0.5$; rozmiar okna z $K(w)=w(w{-}1)/2\le C_{\mathrm{eff}}=180$
daje $w_{\max}=19$; okna sekwencyjnie, scalanie zachłanne po
$|\delta_k|$; umożliwia $n\le50$.

\begin{table}[h]\centering\small
\caption{Traktowalność (skrót): $K$ po filtrze blokowym, największa
wykonalna instancja.}
\begin{tabular}{@{}llll@{}}\toprule
Klasa & $K$ & $n_{\max}$ & Sprzęt\\\midrule
Klasyczne & $O(m)$--$O(nm)$ & $\infty$ & CPU\\
Kwantowe oryg. & $n{-}1$ (adj/fib), $O(n^2)$ (dyn/motz) & 20 & QPU\\
Enhanced & $n{-}1$; $\le171$ okienkowane & 200/500; 50 & QPU\\
\bottomrule\end{tabular}\end{table}

\section{Formulacje QUBO}
Forma ogólna $\min x^\top Qx$, $x\in\{0,1\}^K$, z przejściem do Isinga
$H_C=\sum h_iZ_i+\sum J_{ij}Z_iZ_j$ (Lucas 2014). Diagonala
$Q_{kk}=\delta_k$; kary $Q_{kl}=\lambda=(\sum_k|\delta_k|)+1$.
Adjacent enh.: one-hot ($Q_{ii}=\delta_i$, $Q_{ij}=2\lambda$) ---
dokładnie jeden swap. Fibonacci enh.: QUBO tridiagonalny
($Q_{i,i+1}=\lambda$); rozwiązań $F_{n+1}$. Dynasearch enh.:
$Q_{kl}=\lambda$ dla par nachodzących; nienachodzące niezależne ---
multi-swapy. Motzkin enh.: $Q_{kl}=\lambda$ dla konfliktów wg reguł
Motzkina; rozwiązań $M_n$.

\section{Setup eksperymentalny}
Wszystko stałe poza sąsiedztwem. Benchmarki: Taillard,
$n\in\{20,50,100,200,500\}$, $m\in\{5,10,20\}$, 10 instancji, 5 seedów,
budżety 100--10\,000 ms (klasyka wszystkie sześć; QPU: tl=1000 i 5000
w tabelach, konwergencja 100--5000 dla enhanced). Sprzęt: Apple M4 Pro
(klasyka); D-Wave Advantage\_system4.1 (Pegasus P16, 5627 kubitów,
40\,279 sprzęgieł), EmbeddingComposite, num\_reads=100 na okno.
Praktyczna pojemność gęstych grafów ok.\ 190--200 zmiennych logicznych
(łańcuchy $\sim$30 kubitów fizycznych); granica miękka: $K_{199}$
adjacenta enhanced weszła w 50/50 runów, $\sim$190-zmiennowy dynasearch
padał sporadycznie, a porażki przechodziły przy retry; łańcuchowy
fibonacci embeduje się trywialnie do $n\le500$. Metryka:
$\mathrm{RPD}=(C_{\max}-C^*)/C^*\times100\%$; istotność: Wilcoxon
($\alpha=0.05$) dla par wariantów.

\section{Wyniki i dyskusja}
Każda liczba kwantowa pochodzi z realnych runów Advantage\_system4.1 ---
niczego nie symulujemy.

\begin{table}[h]\centering\small
\caption{Wszystkie 12 sąsiedztw, tai20$\times$5, $t_{\max}=1000$ ms
(średni RPD \%). Kwanty: reżim 1 iteracji (latencja $>$ budżet);
dynasearch 48/50 runów (2 porażki embeddingu); enhanced: 3 okna na
wywołanie, 0 porażek w 400 runach.}
\begin{tabular}{@{}llrr@{}}\toprule
Klasa & Sąsiedztwo & ILS & SA\\\midrule
\multirow{4}{*}{Klasyczne} & Adjacent & 3.89 & 11.28\\
& Fibonacci & 4.09 & 15.00\\ & Dynasearch & 3.32 & 4.19\\
& Motzkin & 3.19 & 4.00\\\midrule
\multirow{4}{*}{Quantum QUBO} & Adjacent & 25.07 & 25.04\\
& Fibonacci & 23.79 & 23.46\\ & Dynasearch & 26.93 & 26.35\\
& Motzkin & 21.66 & 21.50\\\midrule
\multirow{4}{*}{Enhanced} & Adjacent & 25.04 & 25.04\\
& Fibonacci & 23.60 & 23.66\\ & Dynasearch & 25.56 & 25.71\\
& Motzkin & 25.98 & 25.45\\\bottomrule\end{tabular}\end{table}

\paragraph{Baseline klasyczny ($t_{\max}=10$ s, $n=20$).}
Dynasearch i motzkin prowadzą przy każdym $m$, przewaga rośnie z $m$
(0.7--1.2 pp przy $m=5$; 1.2--1.4 przy $m=20$): ILS/SA ---
adj 3.41/6.98 ($m{=}5$), 11.22/15.61 ($m{=}10$), 22.58/26.27 ($m{=}20$);
fib 3.82/9.24, 11.23/17.81, 22.17/26.82; dyn 2.68/3.36, 9.99/10.75,
21.41/21.99; motz 2.64/3.25, 10.21/10.68, 21.46/21.73. SA adj/fib tracą
7--10 pp do ILS przy $m=5$ (oraculum best-move nad rzadkim sąsiedztwem
więzi SA w wąskich basenach); SA dyn/motz w granicach 0.6--0.7 pp.

\paragraph{Kwanty na QPU przy $n=20$.} Każdy run rejestruje liczbę
iteracji --- reżim porównania jest jawny. Przy tl=1000 dokładnie
1 iteracja (wywołanie $\sim$2.1 s ściennie dla adj/fib). Ranking odwraca
się względem klasyki: motzkin 21.7\% $<$ fibonacci 23.8 $<$ adjacent
25.1 $<$ dynasearch 26.9 --- gdy budżet starcza na jeden ruch, wygrywa
najbogatszy zbiór ruchów. Jakość podąża za iteracjami: przy tl=5000
adj/fib robią 3 iteracje, fibonacci 19.9\%, motzkin 19.5\%. Dynasearch
obnaża granicę embeddingu (4/200 runów bez embeddingu przy $\sim$190
gęstych zmiennych; motzkin przy podobnej liczbie, ale rzadszym grafie,
embedował się zawsze) --- wykonalność rządzi się gęstością grafu
konfliktów, nie liczbą zmiennych.

\paragraph{Anatomia wywołania QPU} (1332 wywołania): rozliczany dostęp
$\sim$30 ms (16 programowanie + 14 próbkowanie); ściennie: adj 2.12 s
(udział QPU 1.4\%), fib 2.06 s (1.5\%), motz 18.5 s (0.16\%; budowa QUBO
$O(mn^3)$), dyn 140 s (0.02\%; minorminer). Porównanie ścienne ignorujące
te proporcje mierzy chmurę, nie procesor.

\paragraph{Wiersze enhanced (efekt filtra delty).} Liniowe sąsiedztwa:
enhanced $=$ oryginał (adj 25.04 vs 25.05, fib 23.60/23.66 vs
23.79/23.46; Wilcoxon $p\ge0.25$) --- replikacja między kampaniami.
Okienkowany dynasearch \emph{zyskuje} 0.9 pp (25.78 vs 26.64, $p=0.012$)
i eliminuje porażki embeddingu (0/400 vs 4\%). Okienkowany motzkin
\emph{traci} 4.1 pp (25.72 vs 21.58, $p<10^{-4}$). Odczyt strukturalny:
konflikty dynasearcha są lokalne, granica okna mało odcina; siła
motzkina to zagnieżdżone systemy łuków rozpinające permutację --- okna
tną właśnie je. Przy tl=5000 dynasearch trzyma przewagę (25.07 vs 26.15,
$p=0.005$), luka motzkina rośnie do 5.6 pp (25.06 vs 19.45,
$p<10^{-4}$): strata okienkowania kumuluje się, zamiast amortyzować.

\begin{table}[h]\centering\small
\caption{Skalowanie: średni RPD \%, ILS, $t_{\max}=5000$ ms; $m=10$ dla
$n\le200$, $m=20$ przy $n=500$. $\dagger$ = 1 iteracja (nieporównywalne);
n/a = poza traktowalnością.}
\begin{tabular}{@{}lrrrrr@{}}\toprule
Sąsiedztwo & $n{=}20$ & $n{=}50$ & $n{=}100$ & $n{=}200$ & $n{=}500$\\\midrule
Kl. Adjacent & 11.61 & 10.03 & 9.54 & 10.37 & 15.36\\
Kl. Fibonacci & 11.27 & 7.74 & 6.15 & 6.19 & 11.12\\
Kl. Dynasearch & 10.23 & 5.39 & 4.67 & 11.63 & 14.53$^\dagger$\\
Kl. Motzkin & 10.45 & 13.49 & 14.21 & 13.80 & 16.26$^\dagger$\\\midrule
Q-Adjacent & 34.91 & n/a & n/a & n/a & n/a\\
Q-Fibonacci & 31.60 & n/a & n/a & n/a & n/a\\
Q-Dynasearch & 39.30 & n/a & n/a & n/a & n/a\\
Q-Motzkin & 34.41 & n/a & n/a & n/a & n/a\\\midrule
Q-Adjacent Enh. & 34.86 & 26.70 & 21.73 & 16.40 & n/a\\
Q-Fibonacci Enh. & 31.68 & 23.56 & 19.62 & 14.99 & 15.75\\
Q-Dynasearch Enh. & 38.55 & 26.37 & n/a & n/a & n/a\\
Q-Motzkin Enh. & 38.44 & 26.03 & n/a & n/a & n/a\\
\bottomrule\end{tabular}\end{table}

\paragraph{Skalowanie.} Warianty enhanced budowano, by choć utrzymały
poziom przy rosnącym $n$; pomiary pokazują więcej --- RPD kwantowy spada
monotonicznie. Fibonacci enhanced: $31.7\to23.6\to19.6\to15.0\%$
($n=20\to200$); adjacent enhanced $\sim$1.5 pp za nim ($34.9\to16.4$).
Luka do klasycznego fibonacciego zwęża się: 20.4 pp ($n{=}20$) $\to$
13.5 ($n{=}100$) $\to$ 8.8 ($n{=}200$) $\to$ 4.6 ($n{=}500$), gdzie
wariant kwantowy (15.75\%) osiąga parytet z klasycznym adjacent
(15.36\%). Prawy koniec krzywej umożliwia tridiagonalność: embedding
w milisekundach nawet przy $K_{499}$, trzy iteracje w 5 s przy każdym
rozmiarze; adjacent od $n=100$ spada do 1 iteracji. Kolumna $n=20$
replikuje asymetrię okien przy $m=10$ (dyn: 38.6 vs 39.3, $p=0.18$;
motz: 38.4 vs 34.4, $p<10^{-4}$). Gęsty one-hot skaluje się dalej, niż
sugerowała kampania $n=20$: $K_{199}$ embedował się w 50/50 runów
($\sim$8 s minorminera); porażki $n=20$ były stochastyczne (retry
przechodził). Twarda ściana: $K_{499}$; okolice 200 zmiennych to
loteria o korzystnych szansach, nie klif.

\paragraph{Klasyczny crossover.} Przy $n=100$ wygrywa dynasearch
(4.67 vs 6.15); przy $n=200$ odwrócenie (6.19 vs 11.63) --- DP $O(n^2)$
zjada budżet; motzkin ($O(n^3)$) odpada wcześniej. Przy $n=500$
pojedyncze wywołanie DP (dyn $\sim$200 s, motz $\sim$30 s) przekracza
każdy budżet --- 1 iteracja, wartości dla kompletności. Mimo to jedno
wywołanie dynasearcha bije adjacenta z tysiącami iteracji (14.53 vs
15.36): struktura ruchu znaczy więcej niż liczba iteracji, o ile ruch
niesie dość swapów. Okienkowane dyn/motz kończą na $n=50$ (26.37/26.03,
$\sim$3 pp za fibonaccim): bariera ścienna, nie embedding. Poza $n=200$
fibonacci enhanced jest jedynym kwantowym sąsiedztwem bez obcięć.

\paragraph{Per budżet ($n=20$, $m=5$, bez poolingu).} Klasyka ILS:
6.60/5.96/3.92/4.33 (100 ms) $\to$ 3.41/3.82/2.68/2.64 (10 s). Wiersze
enhanced płaskie 100$\to$1000 ms (1 iteracja; budżet poniżej $\sim$2 s
kosztu drugiego wywołania nic nie kupuje), spadek przy 5000 ms
(2--3 iteracje). Kreski = warstwy jeszcze niezmierzone w protokole transz.

\paragraph{Jakość okienkowania.} Pary przecinające granice okien nigdy
nie są ewaluowane; strata zmierzona przy $n=20$: $\approx$0 dla
dynasearcha, $\sim$4 pp narastająco dla motzkina. Wszystkie runy
z $\rho=0.5$; czułość na $\rho$ niezbadana.

\paragraph{Wkład BTJL.} Ablacja 3600 parowanych runów (oba ramiona,
$n=20$, $m\in\{5,10,20\}$, $t\in\{1,5,10\}$ s): ILS z BTJL bije restart
losowy w każdej komórce o 1.1--2.6 pp ($p<10^{-17}$); rzadkie sąsiedztwa
zyskują 4.1 pp (adjacent) i 3.5 pp (fibonacci), dyn/motz istotne
0.4/0.3 pp ($p<10^{-19}$).

\paragraph{Przepustowość.} Wywołanie okienkowane mnoży 30 ms dostępu
przez liczbę okien; koszt ścienny pozostaje zdominowany przez embedding
i kolejkę.

\section{Wnioski}
Klasa enhanced łamie barierę $n=20$ dwiema zmianami: usunięciem filtra
delty i dekompozycją okienkową. Akcelerator Smutnickiego tnie zmienne
o 30--60\%. BTJL bije restart losowy w każdej komórce ablacji
(1.1--2.6 pp, $p<10^{-17}$). Kampanie QPU z lipca 2026 wypełniły każdą
traktowalną komórkę tabel realnymi pomiarami i odpowiedziały na obie
hipotezy: (i) fibonacci enhanced trzyma poziom w całym zakresie ---
RPD spada monotonicznie $31.7\to15.0\%$, luka do klasyki zwęża się
$20.4\to4.6$ pp przy $n=500$ (parytet z klasycznym adjacent);
(ii) koszt okienkowania jest strukturalny: dynasearch replikuje lub
lekko bije pełny QUBO eliminując porażki embeddingu, motzkin płaci
$\sim$4 pp rosnąco z budżetem (zagnieżdżenia międzyokienne). Przy
$n=50$, gdzie żaden pełny QUBO $O(n^2)$ się nie embeduje, oba
okienkowane warianty dostarczają działające ruchy QPU w granicach 3 pp
od fibonacciego. Reżim jednej iteracji, anatomia kosztu QPU
i stochastyczność embeddingu raportowane są obok jakości, bo definiują,
co uczciwe porównanie ścienne może twierdzić.
\textbf{Future work:} adaptacyjny dobór okien, obwody QAOA na
procesorach bramkowych, dekompozycje zachowujące interakcje
międzyokienne, hybrydowy flow shop.

\end{document}
